% Stages A and B: the specification, the brute-force reference, the command-line program and the
% statements comparator certifies. Proposed names are set with \lname and become \lean once the
% declarations exist.
\chapter{The specification and the brute-force reference}\label{chap:spec}

This chapter states the specification of the points that ratpoints finds, settled in stage~A, and
describes the work of stage~B: the specification in Lean, the brute-force reference and the
command-line program, together with the statements that comparator certifies. The specification is
the one part whose correctness no proof can check, so wherever there is a choice it copies what the
C code computes, and each of its clauses is compared with ratpoints~3.1.1 below
(Section~\ref{sec:spec-evidence} summarizes the evidence). Where the library function
\verb|find_points| and the command-line program differ, as in the order of the checks on the input,
the specification and the model follow the program.

Throughout, $c = [c_0, c_1, \dots, c_n]$ is a list of integers, the coefficient list of the
program's input with the constant term first, so that $n + 1$ is its length; $d$ is its
\emph{degree}, the largest index with $c_d \ne 0$ ($0$ if there is none); and
\[ D = d + (d \bmod 2), \qquad N = 2 \lceil n/2 \rceil = 2 \lfloor (n+1)/2 \rfloor \]
are $d$ and $n$ rounded up to even numbers ($N = 0$ for the empty list). $H$ is a natural number,
the height bound.

\section{The binary form}

\begin{definition}[The binary form]\label{def:binaryForm}
  \lean{Ratpoints.degree, Ratpoints.binaryForm}\leanok
  For integers $x$ and $z$,
  \[ F_c(x, z) = \sum_{i=0}^{d} c_i\, x^i z^{D-i}. \]
\end{definition}

\begin{verbatim}
/-- The degree of the polynomial `∑ cᵢ Xⁱ`: the largest index of a nonzero
entry of `c`, or `0` if there is none. -/
def degree (c : List ℤ) : ℕ := (c.reverse.dropWhile (· == 0)).length - 1

/-- The binary form `F(x, z) = ∑ cᵢ xⁱ z^(D - i)` of even degree
`D = d + d % 2` of the polynomial of degree `d` with coefficient list `c`. -/
def binaryForm (c : List ℤ) (x z : ℤ) : ℤ :=
  ∑ i ∈ Finset.range (degree c + 1),
    c.getD i 0 * x ^ i * z ^ (degree c + degree c % 2 - i)
\end{verbatim}

\begin{remark}[What the C code computes]\label{rem:degree}
  This is the form the C code evaluates. The library drops the leading zero coefficients
  (\csource{find_points.c}{4624}{4632}), so it works with the degree $d$, and its exact check
  computes $\sum_{k \le d} c_k a^k b^{d-k}$ by Horner's rule in the numerator $a$, with the
  products $c_k b^{d-k}$ computed once per denominator $b$, and multiplies by $b$ once more when
  $d$ is odd (\csource{find_points.c}{5291}{5298}); the result is $F_c(a, b)$. For instance
  \verb|'-2 0 0 1' 130 -q| prints $(129 : 3830 : 100)$, and $3830^2 = 100 \cdot (129^3 - 2 \cdot
  100^3) = F_c(129, 100)$. The library refuses every input with
  $\lceil d/2 \rceil < \lceil n/2 \rceil$, so on an accepted input $D = N$; this held in every
  accepted run of the empirical checks, with and without \verb|-s| and \verb|-k|.
\end{remark}

The model gets its own evaluations of the form, the brute-force reference of stage~B and the exact
check of stage~C1, each in the manner of the C code and each proved equal to $F_c$.

\begin{definition}[The binary form of degree $N$]\label{def:binaryFormPoly}
  \uses{def:binaryForm}
  \lean{Ratpoints.binaryFormPoly, Ratpoints.eval_binaryFormPoly}\leanok
  \[ \Phi_c = \sum_{i=0}^{n} c_i X^i Z^{N-i} \in \QQ[X, Z], \]
  the polynomial made homogeneous of degree $N$, which is the form the refusals speak about (the
  C code's ``polynomial not squarefree as a binary form of even degree''). When
  $\lceil d/2 \rceil = \lceil n/2 \rceil$, that is on every input the library accepts and whenever
  $\Phi_c$ is squarefree, $D = N$ and $\Phi_c(x, z) = F_c(x, z)$ for all integers $x$ and $z$.
\end{definition}

\begin{verbatim}
noncomputable def binaryFormPoly (c : List ℤ) : MvPolynomial (Fin 2) ℚ :=
  ∑ i ∈ Finset.range c.length,
    MvPolynomial.C (c.getD i 0 : ℚ) * MvPolynomial.X 0 ^ i *
      MvPolynomial.X 1 ^ (2 * (c.length / 2) - i)

theorem eval_binaryFormPoly (h : c.length / 2 ≤ (degree c + 1) / 2) (x z : ℤ) :
    MvPolynomial.eval ![(x : ℚ), z] (binaryFormPoly c) = (binaryForm c x z : ℚ)
\end{verbatim}

The refusal theorem uses Mathlib's \verb|Squarefree|: no square of a non-unit divides $\Phi_c$. The
units of $\QQ[X, Z]$ are the nonzero constants, so this is squarefreeness over $\QQ$ (over $\ZZ$,
$4(X^2 + Z^2)$ would not be squarefree, while the program accepts \verb|4 0 4|), and $0$ is not
squarefree. The coefficients are written \verb|c.getD i 0|, since Mathlib's notation \verb|c[...]|
for cycles makes \verb|c[i]| fail for a list named \verb|c| wherever all of Mathlib is imported.

$\Phi_c$ is also Mathlib's homogenization \verb|Polynomial.homogenize f N| of
$f = \sum c_i X^i \in \QQ[X]$, which stage~D will use. Mathlib says nothing yet about the
squarefreeness of a homogenization. Stage~D will prove: for a polynomial $p$ over a field and
$m \ge \deg p$, the homogenization of $p$ of degree $m$ is squarefree if and only if $p$ is
squarefree and $m \le \deg p + 1$; together with the facts this needs that Mathlib lacks: a factor
of a nonzero homogeneous polynomial over a domain is homogeneous, \verb|Squarefree| is preserved
by multiplicative equivalences, and the reversal of $f$ swaps $X$ and $Z$ in its homogenization.
These are candidates for \verb|Ratpoints/Mathlib|. With them, ``$\Phi_c$ squarefree'' is
``$N \le d + 1$ and $f$ squarefree'', the two conditions the C code tests.

\section{The points}

\begin{definition}[Normalized point of height at most $H$]\label{def:IsPoint}
  \uses{def:binaryForm}
  \lean{Ratpoints.IsPoint}\leanok
  A triple $(x, y, z)$ of integers is a
  \emph{normalized point of height at most $H$} on $y^2 = F_c(x, z)$ if
  \[ y^2 = F_c(x, z), \quad \gcd(x, z) = 1, \quad z > 0 \text{ or } (z = 0 \text{ and } x = 1),
     \quad |x| \le H, \quad z \le H. \]
\end{definition}

\begin{verbatim}
def IsPoint (c : List ℤ) (H : ℕ) (p : ℤ × ℤ × ℤ) : Prop :=
  p.2.1 ^ 2 = binaryForm c p.1 p.2.2 ∧ Int.gcd p.1 p.2.2 = 1 ∧
    (0 < p.2.2 ∨ p.2.2 = 0 ∧ p.1 = 1) ∧ |p.1| ≤ H ∧ p.2.2 ≤ H
\end{verbatim}

Its clauses, compared with the program:
\begin{itemize}
  \item \emph{The normalization} is the one the program prints. Since $D$ is even,
    $F_c(-x, -z) = F_c(x, z)$, and in the weighted projective plane with weights $(1, D/2, 1)$,
    where the points of the curve live, $(-x : (-1)^{D/2} y : -z)$ is the same point as
    $(x : y : z)$. As both signs of $y$ are listed, the condition can pick $z > 0$, and $x = 1$
    when $z = 0$ (where $\gcd(x, 0) = |x|$ forces $x = \pm 1$).
  \item \emph{The points at infinity} are therefore the triples $(1, y, 0)$ with
    $y^2 = F_c(1, 0) = c_D$: for odd $d$, $c_D = c_{d+1} = 0$ and there is the one point
    $(1 : 0 : 0)$; for even $d$, $c_D = c_d$, which gives $(1 : \pm\sqrt{c_d} : 0)$ when $c_d$ is a
    square, and nothing otherwise. These are exactly the points at infinity of the C code
    (\csource{find_points.c}{5068}{5105}).
  \item \emph{Both signs of $y$}: a point with $y \ne 0$ gives the two triples $(x, \pm y, z)$,
    as the program prints both.
  \item \emph{The box is closed}: $|x| = H$ and $z = H$ are included, as in the C code (the
    numerators run over $-H, \dots, H$ and the denominators over $1, \dots, H$).
  \item \emph{Coprimality} is \verb|Int.gcd|, with $\gcd(0, z) = z$, so $x = 0$ occurs only as
    $(0, y, 1)$, as in the C code's test (\csource{sift.c}{318}{343}).
  \item \emph{The height bound applies to the points at infinity too} ($|1| \le H$), so for
    $H = 0$ there are no points at all. The program refuses $H = 0$, so this never shows; it
    keeps the predicate uniform.
\end{itemize}

\begin{definition}[The set of points]\label{def:points}
  \uses{def:IsPoint}
  \lean{Ratpoints.points, Ratpoints.finite_setOf_isPoint, Ratpoints.mem_points}\leanok
  The finite set of normalized points of height at most $H$ on
  $y^2 = F_c(x, z)$. It is finite because $|x| \le H$, $0 \le z \le H$ and
  $|y| \le y^2 = F_c(x, z) \le \sum_i |c_i| H^D$.
\end{definition}

\begin{verbatim}
theorem finite_setOf_isPoint (c : List ℤ) (H : ℕ) : {p | IsPoint c H p}.Finite

noncomputable def points (c : List ℤ) (H : ℕ) : Finset (ℤ × ℤ × ℤ) :=
  (finite_setOf_isPoint c H).toFinset

theorem mem_points : p ∈ points c H ↔ IsPoint c H p
\end{verbatim}

\verb|IsPoint| is decidable (through a named instance built with \verb|inferInstanceAs|, since
instance search does not unfold a definition), and \verb|degree| and \verb|binaryForm| are
computable, so the specification file can check facts taken from the program's output by
\verb|decide|: for instance that $(129, 3830, 100)$ is a point of $y^2 = x^3 - 2$, that
\verb|1 1 1 1 0| (degree $3$, the one vanishing leading coefficient allowed for even $n$) has the
point $(1, 0, 0)$, and that $(1, 1, 0)$ is not a point of \verb|1 0 0 0 2|, whose leading
coefficient $2$ is not a square.

\section{The refusals}\label{sec:refusals}

The refusal theorem belongs to stage~D, but its predicate is part of the specification:
\begin{verbatim}
theorem exists_findPoints_eq_error_iff (hP : P.WF) :
    (∃ e, findPoints c H P = .error e) ↔
      c.length < 2 ∨ 101 < c.length ∨ H = 0 ∨ ¬ Squarefree (binaryFormPoly c)
\end{verbatim}
In a default run with $2$ to $101$ coefficients and $1 \le H \le 2^{63} - 1$, the program refuses
exactly the inputs for which $\Phi_c$ is not squarefree. The reason: $\Phi_c = Z^{N - d}\,
\Phi^0$, where $\Phi^0$ is the homogenization of $f = \sum c_i x^i$ in degree $d$, which is
prime to $Z$; so $\Phi_c$ is squarefree if and only if $N - d \le 1$ and $f$ is squarefree. The
library tests the first condition when it drops the leading zeros
(\csource{find_points.c}{4624}{4632}) and the second with the remainder sequence of $f$ and $f'$
in the Sturm step (\csource{sturm.c}{120}{143}); its test of $c_0 = c_1 = 0$ in the reversal
(\csource{find_points.c}{4744}{4749}) catches $X^2 \mid \Phi_c$, which the remainder sequence
would also catch.

The model keeps the bound of 101 coefficients (\verb|RATPOINTS_MAX_DEGREE| $= 100$), which 3.1.1
imposes in the library (\csource{find_points.c}{4557}{4558}) as well as in the program. In the
order of the checks the library and the program differ, and the model follows the program: the
library checks the leading zeros before the height, so that \verb|5 0| with $H = 0$ is
\verb|NON_SQUAREFREE| for it, while
the program checks the coefficient list and the height before it calls the library
(\csource{main.c}{362}{380}) and reports error~4. So \verb|findPoints| refuses more than 101
coefficients and fewer than two, then $H = 0$, and then a form that is not squarefree, with the
program's error numbers 3, 5, 4 and~8.

With the options of stage~D the predicate changes: under \verb|-s| (no Sturm step) the program
refuses, besides the coefficient count and $H = 0$, exactly the inputs with $N - d \ge 2$ and,
when the reversal check runs (no \verb|-k|, no search intervals, the full range of denominators),
those with even $d$ and $c_0 = c_1 = 0$; under \verb|-s -k| only those with $N - d \ge 2$.

\section{The brute-force reference}\label{sec:bruteForce}

\begin{definition}[Brute-force search]\label{def:bruteForce}
  \lean{Ratpoints.bruteForce}\leanok
  The list, computed in the model, of the triples found by testing, in this order, the candidate
  $(1, 0)$ if $H \ge 1$, and then the pairs $(x, z)$ with $1 \le z \le H$, $-H \le x \le H$ and
  $\gcd(x, z) = 1$, with $z$ ascending and, for each $z$, $x$ ascending. For each candidate, $v =
  F_c(x, z)$ is computed; if $v \ge 0$ and $v = y^2$ for $y = \lfloor \sqrt{v} \rfloor$, the triple
  $(x, y, z)$ is output, followed by $(x, -y, z)$ if $y \ne 0$.
\end{definition}

This is the core of \verb|verify-test4.py| of the C repository (which compares point sets
only), in the model (Lean core, no Mathlib). It evaluates $F_c$ as the C code's exact check does
(Horner's rule in $x$ on the products $c_k z^{d-k}$, computed once per $z$, and one more factor
$z$ when $d$ is odd), so that the candidate $(1, 0)$ needs no case of its own, and it takes the
square root with \verb|Nat.sqrt|. It takes no hypothesis on $c$ or $H$, and it models no C code,
so it cites none. Written as nested \verb|List.flatMap|s rather than loops, its membership proof
is direct. A compiled prototype tests about $10^6$ coprime pairs per second (Lean's \verb|Int| is
a GMP number beyond $\pm 2^{31}$, which $F_c$ exceeds at small heights already), so whole test
suites fit at heights up to about $1000$; there are about $1.2\, H^2$ coprime pairs per curve.
Its order is the program's for the inputs that the program neither reverses nor searches in
divisor order: on the \verb|test4.sh| runs of height at most $1000$ that use only the
coefficients, the height, \verb|-q| and tuning options, it reproduces the C output byte for byte
in 167 of 195 cases, and the other 28 (reversed polynomials, and odd degree with $|c_d| \ne 1$)
agree as sets. The order of the program is a theorem of stage~G about \verb|findPoints|; the
reference is compared as a set.

\begin{theorem}[The brute-force reference finds exactly the points]\label{thm:bruteForce}
  \uses{def:bruteForce, def:IsPoint, def:points}
  \lean{Ratpoints.mem_bruteForce_iff, Ratpoints.nodup_bruteForce, Ratpoints.toFinset_bruteForce}
  \leanok
  For every list $c$ and every $H$: a triple lies in
  $\mathrm{bruteForce}(c, H)$ if and only if it is a normalized point of height at most $H$;
  the list has no duplicates; and its set of elements is $\mathrm{points}(c, H)$.
\end{theorem}

\begin{proof}
  \leanok
  The candidate $(1, 0)$ gives the points with $z = 0$, since $\gcd(x, 0) = 1$ and the
  normalization force $x = 1$; the pairs give those with $z > 0$. The model's evaluation is
  $F_c$, and $v = (\lfloor \sqrt{v} \rfloor)^2$ holds exactly when $v \ge 0$ is a square. The
  candidates are distinct, and $y \ne -y$ for $y \ne 0$, so there are no duplicates. The third
  statement follows from the first two.
\end{proof}

\begin{verbatim}
theorem mem_bruteForce_iff : p ∈ bruteForce c H ↔ IsPoint c H p
theorem nodup_bruteForce (c : List ℤ) (H : ℕ) : (bruteForce c H).Nodup
theorem toFinset_bruteForce (c : List ℤ) (H : ℕ) :
    (bruteForce c H).toFinset = points c H
\end{verbatim}

\section{What comparator certifies}\label{sec:comparator-ab}

\verb|comparator/Challenge.lean| imports Mathlib and the model. It restates \verb|degree|,
\verb|binaryForm| and \verb|IsPoint| verbatim and states, with \verb|sorry|,
\begin{verbatim}
theorem challenge_mem_bruteForce_iff (c : List ℤ) (H : ℕ) (p : ℤ × ℤ × ℤ) :
    p ∈ Ratpoints.bruteForce c H ↔ Ratpoints.IsPoint c H p
theorem challenge_nodup_bruteForce (c : List ℤ) (H : ℕ) :
    (Ratpoints.bruteForce c H).Nodup
\end{verbatim}
\verb|comparator/Solution.lean| proves them by the library's theorems. \verb|binaryFormPoly|
joins the challenge with the refusal theorem (stage~D), and \verb|points| stays out of it:
comparator compares the full declarations that a certified statement reaches, the proofs of
theorems included, so the finiteness proof inside \verb|points| would have to be reproduced
term for term; and the two theorems above say everything the \verb|Finset| statement says. (A
computable \verb|points|, filtering an explicit box with the bound $\max |F_c|$ for $|y|$,
would avoid the proof term, but it is not needed.) The specification can be restated because
it is elaborated with Mathlib, as the challenge is; a definition of the model, elaborated
without Mathlib, has other instance terms and is always imported, never restated. Membership and
no duplicates for \verb|findPoints| join in stage~C1, together with the well-formedness of the
default parameters.

\section{The command-line program}\label{sec:cli}

The Lean program is a program of its own, the executable \verb|ratpoints-lean|, with the command
line of ratpoints. Eventually it offers the full functionality of ratpoints except the threads
and the other hardware-dependent features. In stage~B it prints the points of the brute-force
reference, so that the test scripts of the C repository can run it as \verb|RP=ratpoints-lean|:
\begin{itemize}
  \item it behaves as with \verb|-q|, printing exactly one line \verb|(x : y : z)| per point and
    nothing else; it accepts \verb|-q|, and \verb|-v|, which \verb|-q| cancels in \verb|main.c|;
  \item it reads the first two arguments as \verb|main.c| does: the coefficient list as one
    argument and the height, with its checks, messages and exit codes, in its order: fewer than
    two arguments (2, with the banner and the usage), a malformed coefficient (11), more than
    101 coefficients (3), fewer than two (5), a height outside $1, \dots, 2^{63} - 1$ (4)
    (\csource{main.c}{362}{380}, \csource{main.c}{700}{740}, \csource{main.c}{1036}{1056});
    every message goes to standard output as a blank line, the message and a blank line;
  \item it reads the tuning options \verb|-n -N -P -p -F -j -J -r -R -U -A -C -Q -W -t| with the
    validation of \verb|main.c| (error~6 and the usage for a missing or malformed argument of the
    thirteen that take one, \csource{main.c}{437}{532}; \verb|-j| and \verb|-J| take none) and
    ignores them for now (\verb|main.c| itself never reads them again under \verb|-q|; in 373 runs
    of the C program they did not change the \verb|-q| output); from stage~C1 on,
    \verb|-n -N -P -p -F -j -J| go into the sieving parameters \verb|P|, while the others, which
    tune only the C program's heuristics and its threads, stay ignored;
  \item it answers every other option of ratpoints, which it does not support yet, with the
    message \verb|Option <opt> is not supported by ratpoints-lean yet.| and the exit status 12,
    which \verb|main.c| does not use (an ignored search interval or denominator range, \verb|-l|,
    \verb|-u|, \verb|-dl| or \verb|-du|, on one of the runs with heights up to $2^{63} - 1$ would
    otherwise turn into a search of the whole box that never ends), and an
    unknown option with error~6, as \verb|main.c| does;
  \item its banner starts with the line \verb|This is ratpoints-lean-0.1|, so that
    \verb|test4.sh| filters it as it does the C program's, and the harness ignores the rest of
    the banner;
  \item it does not refuse a polynomial that is not squarefree: that is the squarefree test of
    stage~D; the test harness skips the inputs that the C program refuses.
\end{itemize}
The numbers need a reader of their own that mirrors \verb|scan_mpz| and \verb|scan_long|
(\csource{main.c}{700}{740}): \verb|main.c| accepts a leading \verb|+| and any C white space
around and between the numbers (\verb|test4.sh| has a tab-separated list and a height ending in
a carriage return), which core's \verb|String.toInt?| does not; and core's \verb|String.toNat?|
accepts the digit separator \verb|_|, which \verb|main.c| refuses. The program lives in the model
library, in \verb|Ratpoints/Model/Cli.lean|, with a \verb|Main.lean| that imports only the model,
as the \verb|lean_exe| target \verb|ratpoints-lean|. Its parts, with the lines of \verb|main.c|
that they mirror:

\begin{definition}[Reading numbers]\label{def:cliNumbers}
  \lean{Ratpoints.Cli.intOfWord?, Ratpoints.Cli.onlyWord?, Ratpoints.Cli.longOfString?, Ratpoints.Cli.isDouble}
  \leanok
  The readers of the numbers of the command line, on the words of an argument, its maximal runs
  of characters other than C white space (the space and the characters from tab to carriage
  return): an integer of the coefficient list as \verb|scan_mpz| reads it, an optional sign
  \verb|+| or \verb|-| and at least one decimal digit, of any size
  (\csource{main.c}{700}{717}); the one word of an argument with only white space around it,
  which \verb|strtol| and \verb|strtod| skip before a number and \verb|only_space_after| checks
  after it (\csource{main.c}{724}{728}); a C \verb|long| as \verb|scan_long| reads it, an
  integer from $-2^{63}$ to $2^{63} - 1$ (\csource{main.c}{733}{740}); and whether an argument
  is a real number as \verb|scan_double| reads it with \verb|strtod| in the C locale: a decimal
  or hexadecimal floating constant or an infinity, a value beyond the range of the doubles
  included, but no NaN (\csource{main.c}{746}{752}).
\end{definition}

\begin{definition}[The failures]\label{def:cliFailure}
  \lean{Ratpoints.Cli.Failure, Ratpoints.Cli.Failure.message, Ratpoints.Cli.Failure.status, Ratpoints.Cli.banner, Ratpoints.Cli.usage}
  \leanok
  The ways in which \verb|ratpoints-lean| refuses its command line, with the message and the
  exit status of each. These are the errors of \verb|main.c| that the command line can cause, 2
  (fewer than two arguments), 3 (more than 101 coefficients), 4 (a height bound that is not a
  C \verb|long| of at least $1$), 5 (fewer than two coefficients), 6 (a malformed option) and
  11 (a coefficient that is not an integer), on which \verb|main.c| prints a blank line, the
  message and a blank line, and exits with the number of the error
  (\csource{main.c}{1036}{1043}, \csource{main.c}{1050}{1053}, \csource{main.c}{1054}{1056});
  on error~2 it prints the banner (\csource{main.c}{43}{48}) and the usage
  (\csource{main.c}{77}{88}) instead of a message, and on error~6 both. The banner is that of
  \verb|ratpoints-lean| and the usage that of ratpoints, with the name of the program changed.
  The last failure is an option of ratpoints that \verb|ratpoints-lean| does not support yet,
  with the exit status 12.
\end{definition}

\begin{definition}[The command line]\label{def:cliInput}
  \uses{def:cliNumbers, def:cliFailure}
  \lean{Ratpoints.Cli.coefficients, Ratpoints.Cli.heightBound, Ratpoints.Cli.checkOptions, Ratpoints.Cli.input}
  \leanok
  The coefficient list and the height bound of a command line, or its first failure, in the
  order of \verb|main.c|: fewer than two arguments fail (\csource{main.c}{210}{214}); the first
  argument gives the coefficients, constant term first, integers separated by white space, and
  fails on a word that is not an integer, wherever it stands, then on more than 101 integers,
  then on fewer than two (\csource{main.c}{362}{375}); the second gives the height bound, a C
  \verb|long| of at least $1$ (\csource{main.c}{378}{379}); and the options are read from left
  to right, each told by the character after its \verb|-| and, for \verb|-f| and \verb|-d|, by
  the one after that (\csource{main.c}{393}{397}, \csource{main.c}{667}{668}). The tuning
  options, of which \verb|-p -F -t -n -N -R -A -P| take a C \verb|long| and
  \verb|-r -U -C -Q -W| a real number, have their arguments checked and are ignored
  (\csource{main.c}{437}{532}), as are \verb|-q|, \verb|-v|, \verb|-j| and \verb|-J|
  (\csource{main.c}{578}{593}); an option of ratpoints that is not supported yet ends the
  reading, its arguments unchecked, and any other word fails with error~6.
\end{definition}

\begin{definition}[The program]\label{def:cliRun}
  \uses{def:cliInput, def:cliFailure, def:bruteForce}
  \lean{Ratpoints.Cli.pointLine, Ratpoints.Cli.run, Ratpoints.Cli.main}
  \leanok
  \verb|ratpoints-lean| reads its command line. On a failure it prints the failure's message and
  exits with its status; otherwise it prints, as ratpoints does with \verb|-q|, the line
  \verb|(x : y : z)| of the default format of \verb|main.c| (\csource{main.c}{681}{685}) for
  each triple of $\mathrm{bruteForce}(c, H)$, in the order of that list, and exits with
  status~0.
\end{definition}

\section{The differential tests of stage~B}

The script \verb|scripts/difftest.py| compares \verb|ratpoints-lean| with the C program built at
the pin and with its recorded outputs, in six suites:
\begin{itemize}
  \item \verb|test4.sh| against \verb|testbase4|, split at the \verb|#| header lines (five of
    which are glued to the end of the preceding output line): 533 of its 582 runs; the others
    are 21 whose output the script filters, 9 that the C program refuses as not squarefree,
    which wait for stage~D, and 19 with heights from $30000$ to $100000$, which wait for
    stage~C1;
  \item \verb|test3.sh| invocation by invocation against the C program, since it has no announce
    lines: 13 of its 21 runs (5 have filtered outputs, one is a refusal, and two have heights too
    large for the brute force);
  \item the curves of \verb|testdata.h|, \verb|testdata-many.h| and \verb|testdata-degrees.h|,
    read from the pinned checkout, at a reduced height $h$ ($100$ by default): the C program run
    at height $h$ prints exactly the reference lines of \verb|testbase*| restricted to
    $|a| \le h$ and $b \le h$, in the same order (checked for $h = 10, 100, 1000$), so these 1206
    runs need no run of the C program;
  \item 300 random polynomials of degree 1 to 8 with coefficients from $-10$ to $10$ at heights
    up to 30, against the C program (from stage~C1 on, when \verb|ratpoints-lean| runs the sieve,
    also the brute force against both at tiny heights);
  \item 219 hand-written command lines (\verb|scripts/difftest_cases.py|) for the checks that
    \verb|main.c| makes on the coefficients, the height and the options: white space and signs,
    the bounds of a \verb|long|, the grammar of \verb|strtod|, every option with a missing or
    malformed argument, and the order of the checks; and the usage of \verb|ratpoints-lean|;
  \item 700 random command lines, with words that the readers of numbers of \verb|main.c| accept
    or refuse in many ways, against the C program.
\end{itemize}
Point lists are compared as sets (the order of the program is a theorem of stage~G about
\verb|findPoints|), and for a run without \verb|-q| with the point lines of the report of the C
program; messages and exit codes are compared exactly, with the banner and the usage of the two
programs exchanged. On a command line with an option that \verb|ratpoints-lean| does not support
yet, it must fail as the C program does if that fails before the option, and else exit with
status~12. CI builds the C program at the pin (\verb|make|, with \verb|libgmp-dev|) and the Lean
program, and runs the script.

\section{Files, names and order of work}

All declarations live in the namespace \verb|Ratpoints|. The module says which part a
declaration belongs to (the model, the specification or the proofs), and
\verb|scripts/check_layout.py| keeps the parts' imports apart. One commit per item, each passing
the project's checks:
\begin{enumerate}
  \item[B1] \verb|Ratpoints/Spec/Basic.lean|, with the definitions and theorems of this chapter's
    first two sections (from \verb|degree| to \verb|mem_points|) and the examples by
    \verb|decide|; the blueprint nodes of those sections become \verb|\lean| nodes.
  \item[B2] \verb|Ratpoints/Model/BruteForce.lean|: \verb|bruteForce|, with \verb|#guard|
    examples taken from the C program.
  \item[B3] \verb|Ratpoints/Proof/BruteForce.lean|: the three theorems of
    Theorem~\ref{thm:bruteForce}, each with an axiom guard (\verb|#print axioms| under
    \verb|#guard_msgs|); the comparator challenge, solution and \verb|challenges.json|;
    \verb|formalization.yaml|.
  \item[B4] \verb|Ratpoints/Model/Cli.lean|, \verb|Main.lean| and the \verb|lean_exe| target
    \verb|ratpoints-lean|, with the correspondence entries for \verb|main.c|.
  \item[B5] \verb|scripts/difftest.py|, with the command lines of
    \verb|scripts/difftest_cases.py|, and its CI job.
\end{enumerate}

\section{The evidence}\label{sec:spec-evidence}

Four independent checks against ratpoints~3.1.1 (tag \verb|v3.1.1|, the pin), run while this
chapter was written:
\begin{itemize}
  \item a literal Python implementation of the specification above compared with the program on
    7964 runs: 133 edge cases and 1200 random polynomials of degree 1 to 8 with heights up to 30,
    each with no option, \verb|-s|, \verb|-k| and \verb|-s -k|; 300 polynomials with a repeated
    factor; 120 curves with large coefficients (up to $2^{128}$, including $-2^{63}$ as the
    leading coefficient); 45 curves that tell the divisor order from ascending order; 300 curves
    at heights 31 to 120. Of these, 6011 runs were accepted and printed 69631 points: no point
    missing, none extra, no duplicate, including 719 accepted runs of polynomials that are not
    squarefree (under \verb|-s| or \verb|-s -k|), which confirms that the output theorem needs no
    squarefreeness. The refusals agreed exactly with the rules of Section~\ref{sec:refusals} (this
    check found that the rule for $c_0 = c_1 = 0$ under \verb|-s| applies to even $d$ only);
    $D = N$ on every accepted run; a model of the program's output order (the points at infinity
    first, then the denominators in the order of the loop, ascending or divisor by divisor, the
    numerators ascending, $+y$ before $-y$, and the effect of a reversal) matched all 6011 runs;
    and the output, order included, did not change under 19 settings of the sieving parameters
    and thread counts (14250 runs);
  \item a second implementation, with an exact model of the output order and Magma deciding
    squarefreeness in $\QQ[x, z]$ independently: the refusal set of a default run is exactly
    ``$\Phi_c$ not squarefree'' on 2337 polynomials (1223 refused), and the order model agreed
    with the program on 696 random runs, including 111 reversed ones;
  \item a compiled Lean prototype of the brute-force reference (core only), which agreed with
    the program as sets on the 195 \verb|test4.sh| runs of Section~\ref{sec:bruteForce}, on the
    1206 curves of the \verb|rptest| suites at heights 10 and 100, and on all their curves at
    height 1000; a model of the C order (reversal, divisor order) reproduced every one of these
    lists exactly;
  \item a reading of \verb|main.c| against the shipped binary, which reproduces \verb|testbase3|
    and \verb|testbase4| exactly, for the command line of Section~\ref{sec:cli}.
\end{itemize}
